% ===================================================================
%  TABELL Nr. 9 — De Bierg. De Bësch.
%  Traduction 2026 d'apres texto/io, controlee sur texto/fr, texto/de
%  et texto/nl. Consigne : voir texto/lb/10-tabell-01.tex.
%
%  L'ECUREUIL DE L'ALINEA 10 EST LE MOT LE PLUS IMPORTANT DE CETTE
%  COLONNE, ET IL FALLAIT NEUF TABLEAUX POUR Y ARRIVER.
%  « Kaweechelchen » : c'est le mot par lequel, pendant l'occupation,
%  on reconnaissait a l'oreille qui n'etait pas d'ici. Un Allemand ne
%  le prononce pas.
%
%  ET C'EST EXACTEMENT LE PROBLEME QUE CETTE COLONNE POSE DEPUIS SON
%  PREMIER TABLEAU, RETOURNE. Le balayage de texto/lb ne peut pas
%  travailler sur le code du caractere, parce que l'allemand et le
%  luxembourgeois s'ecrivent avec le meme alphabet et qu'une phrase
%  allemande passerait sans qu'un signe la denonce. « Kaweechelchen »
%  est le contre-exemple : un mot qui ne trahit rien a l'oeil et
%  tout a l'oreille. CE QUI SEPARE DEUX LANGUES PROCHES NE S'ECRIT
%  PAS TOUJOURS.
%
%  LE BIERG DE LA PREMIERE SCENE N'EST PAS UNE MONTAGNE. Le point le
%  plus haut du pays fait cinq cent soixante metres, et l'Éislek n'a
%  ni glacier, ni avalanche, ni cime. Le luxembourgeois compose donc
%  ce qu'il n'a pas — « Biergkette » la chaine, « Bierghang » le
%  versant, « Biergspëtz » la cime — et emprunte le reste :
%  « Gletscher », « Lawin », « Vulkan », « Krater », « Grott ». C'est
%  l'operation lituanienne du tableau 9, faite sur le meme tableau,
%  par une autre langue et pour la meme raison.
%
%  ET LA SECONDE MOITIE RENVERSE LA PREMIERE, comme en lituanien. Le
%  pays est couvert de bois, et la page devient dense et propre :
%  « Eech » le chene, « Buch » le hetre, « Bierk » le bouleau,
%  « Lann » le tilleul, « Ulm » l'orme, « Moos » la mousse,
%  « Faarn » la fougere, « Heed » la bruyere, « Spiecht » le pic,
%  « Ekster » la pie, « Hues » le lievre, « Réi » le chevreuil,
%  « Wëllschwäin » le sanglier, « Seechomesenhaascht » la
%  fourmiliere.
% ===================================================================

%%K t09-apar-1 apar
\VUsaut{4.0mm}
\VUtitre{15.0pt}{60}{TABELL Nr. 9}
\VUsaut{3.0mm}

%%K t09-tit-1 sub
\VUcentre{12.0pt}{0}{\textit{Éischt Zeen.}}
\VUsaut{3.0mm}

%%K t09-tit-2 sub
\VUpk{11.4pt}{40}{De Bierg. --- D'Buedstad.}
\VUsaut{3.0mm}

%%K t09-c1-01-1 p
1. --- Zënter aacht Deeg sinn ech zu Pratz. Ech kann déi ganz Bewonnerung net
ausdrécken, déi ech gespuert hunn, wéi ech virun der wonnerbarer
\VUgras{Biergkette} \textsuperscript{(1)} vun de Pyrenäe stoung, vun där de
\VUgras{Kamm} \textsuperscript{(2)} um bloen Himmel wéi eng risech Girlande ausgesäit.

%%K t09-c1-02-1 p
2. --- Ech hu mer net virgestallt, datt d'Pyrenäen-\VUgras{Gipfelen}
\textsuperscript{(3)} esou eng grouss \VUgras{Héicht} erreechen, an ech si ganz
erstaunt bliwwe virun de wonnerschéine \VUgras{Gletscher} \textsuperscript{(5)} a
virun de verschneite \VUgras{Biergspëtzten} \textsuperscript{(4)}, op deenen esou
schrecklech \VUgras{Lawinen} erofrullen.

%%K t09-tit-3 sub
\VUsaut{3.0mm}
\VUpk{10.2pt}{40}{Bierglandschaft.}
\VUsaut{3.0mm}

%%K t09-c1-03-1 p
3. --- Direkt den Dag no menger Arrivée sinn ech bei den alen erloschene
\VUgras{Vulkan} \textsuperscript{(6)} gaangen, dee bei Pratz ass. Op de kalen a
plakege Ränner vum \VUgras{Krater} gesäit ee nach gutt d'Spuere vum Duerchgang vun
der \VUgras{Lava,} déi virun e puer Joerhonnerten um \VUgras{Bierghang}
\textsuperscript{(7)} geflass ass. Et ass mer gelongen, op engem vun den
\VUgras{Häng} e puer Stécker vun där Lava ze fannen.

%%K t09-c1-04-1 p
4. --- Pratz läit op der Grenz, an d'\VUgras{Festung} \textsuperscript{(8)}, déi den
\VUgras{Enkpass} \textsuperscript{(27)} verdeedegt, deen Frankräich vu Spuenien
trennt, erënnert ganz un déi al Schlässer um Ufer vum Rhäin, déi vun uewen op hirem
Fiels no \VUgras{Beute} ze lauere schéngen.

%%K t09-c1-05-1 p
5. --- En enormen \VUgras{Adler} \textsuperscript{(9)}, deen säin Nascht net wäit vum
\VUgras{Fort} \textsuperscript{(8)} huet, zitt dacks meng Opmierksamkeet un.

Dee \VUgras{Raubvull,} deem säi \VUgras{Schniewel} staark ass, bréngt senge Jonken
dacks e Lämmche oder e Zicklein, dat hie mat senge \VUgras{Krallen} hält. Sou grouss
sinn seng \VUgras{Flilleken,} datt hie laang mat senger schwéierer Laascht gläide
kann.

%%K t09-tit-4 sub
\VUsaut{3.0mm}
\VUpk{10.2pt}{40}{Den Ausflügler.}
\VUsaut{3.0mm}

%%K t09-c1-06-1 p
6. --- Do ass den agreabelste Spadséiergang den Ausfluch bei d'Kaskad
\textsuperscript{(10)} vu „Kau“. De \VUgras{Waasserfall} fält wirbelnd erof a
verschwënnt duerno als schéine \VUgras{Baach} an engem \VUgras{Dännebësch}
\textsuperscript{(11)}, dee westlech vun der Stad läit. Mir sinn duerch dee Bësch bis
bei d'\VUgras{meteorologesch Observatoire} \textsuperscript{(12)} geklommen, a vun
uewen um \VUgras{Belvedere} hu mir dat ganzt \VUgras{Panorama} vun der Géigend gesinn.
D'\VUgras{Landschaft} ass wonnerbar, an d'\VUgras{Natur} schéngt dem Land all
d'Schätz ze schenken, iwwer déi si verfüügt.

%%K t09-c1-07-1 p
7. --- Fir erofzekommen, hu mir de \VUgras{Funiculaire} \textsuperscript{(13)}
geholl, a mir hunn an enger klenger \VUgras{Statioun} nieft de \VUgras{Ruinen}
\textsuperscript{(14)} vun engem alen \VUgras{Feudalschlass} gehalen, dat fréier vun
den Häre vu Pratz bewunnt gouf.

%%K t09-c1-08-1 p
8. --- Aarmt Schlass! Wat denkt et sech wuel, wann et do ënnen déi koketten
\VUgras{Buedstad} \textsuperscript{(15)} gesäit, déi elo eng modesch
\VUgras{Thermalstatioun} ass?

%%K t09-c1-08-2 p
D'\VUgras{Klima} ass esou agreabel, dat \VUgras{Mineralwaasser} esou wierksam, datt
d'\VUgras{Kur,} déi am \VUgras{Sanatorium} \textsuperscript{(17)} gemaach gëtt, e
richtegt Vergnügen ass.

%%K t09-tit-5 sub
\VUsaut{3.0mm}
\VUpk{10.2pt}{40}{Agreabel Thermalstad.}
\VUsaut{3.0mm}

%%K t09-c1-09-1 p
9. --- Iwwregens huet ee alles gemaach, fir den Openthalt agreabel ze maachen. E
grousse \VUgras{Viadukt} \textsuperscript{(16)} gouf gebaut, fir eng Eisebunn
méiglech ze maachen; de \VUgras{Park} \textsuperscript{(18)} vun der
\VUgras{Thermalanlag,} an deem \VUgras{concertéiert} gëtt, ass charmant ugeluecht. E
puer zierlech „\VUgras{Villaen}“ \textsuperscript{(19)} goufe ronderëm de Casino
\textsuperscript{(21)} gebaut, an eng ganz gutt ënnerhalen \VUgras{Chaussée}
\textsuperscript{(22)} mécht d'Verbindungen immens vill méi einfach.

%%K t09-c1-10-1 p
10. --- Eng einfach \VUgras{Béischung} \textsuperscript{(23)} trennt de
\VUgras{Wee} \textsuperscript{(20)} vum eigentlechen \VUgras{Dall}
\textsuperscript{(25)}, an deem sengem Grond e zierleche klenge \VUgras{Séi}
\textsuperscript{(26)} läit, deen e kloere \VUgras{Baach} \textsuperscript{(24)}
späist.

%%K t09-c1-11-1 p
11. --- Op der anerer Säit vum Dall gëtt eng \VUgras{Eisenminn} \textsuperscript{(28)}
ausgebeut. Ech sinn e puer Mol d'\VUgras{Miner} kucke gaangen, wéi si an de
\VUgras{Schacht} erofgaange sinn, an ech sinn esouguer an d'\VUgras{Galerie}
gaangen, an där si hiren Dag verbréngen, andeems si d'\VUgras{Äerz} erausschaffen,
an deem d'\VUgras{Metall} ass.

%%K t09-c1-12-1 p
12. --- Eng wichteg \VUgras{Gisserei} gouf nieft der Minn gebaut, fir d'Räichtum,
déi doraus geholl gëtt, direkt ze notzen. Den \VUgras{Héichuewen}
\textsuperscript{(29)}, dee mat de \VUgras{Steekuelen} gehëtzt gëtt, déi hei am
Iwwerfloss sinn, mécht et méiglech, séier a bëlleg \VUgras{Gusseisen} ze kréien.

%%K t09-c1-13-1 p
13. --- Net wäit vun der Minn gëtt e \VUgras{Steebroch} \textsuperscript{(30)}
ausgegruewen. Weder \VUgras{Granit,} nach \VUgras{Porphyr,} nach esouguer
\VUgras{Sandsteen} schleet den ale \VUgras{Steebriecher} mat sengem \VUgras{Pickel}
of, mä einfach \VUgras{Quadersteng} fir de Bau vun Haiser.

%%K t09-c1-14-1 p
14. --- Ee vun de schéinsten Zierden vun deem Land sinn d'\VUgras{Dännebeem}
\textsuperscript{(31)}. Iwwerall am Grond vun enger \VUgras{Schlucht}
\textsuperscript{(32)}, um Ufer vun engem \VUgras{Wëldbaach} \textsuperscript{(33)},
vun engem \VUgras{Ofgrond} \textsuperscript{(34)} oder vun engem Precipice ginn hir
dënn Nolen dem Land eng gréng Faarf.

%%K t09-tit-6 sub
\VUsaut{3.0mm}
\VUpk{10.2pt}{40}{Zu Fouss goen.}
\VUsaut{3.0mm}

%%K t09-c1-15-1 p
15. --- Ech notze meng Vakanz, fir \VUgras{laang} ze spadséieren. D'\VUgras{Weeër}
\textsuperscript{(35)}, déi sech de Bierg erop wannen, maachen et méiglech, ouni op
d'Distanz ze kucken, e puer \VUgras{Kilometer} an dacks e puer \VUgras{véierfach
Kilometer} ze goen.

%%K t09-c1-15-2 p
\VUgras{Grenzsteng} \textsuperscript{(36)} a \VUgras{Weeweiser} \textsuperscript{(37)}
markéieren d'Weeër, op deenen ee dacks e jonke \VUgras{Biergbewunner}
\textsuperscript{(38)} begéint, mat engem \VUgras{Rucksak} \textsuperscript{(40)} un
der Säit, dee e \VUgras{Murmeldéier} \textsuperscript{(39)} dréit, oder e
\VUgras{Zigeiner} \textsuperscript{(41)}, deem säi \VUgras{Pelzjackett}
\textsuperscript{(42)} komesch mat der aler zerrasser \VUgras{Kapuz}
\textsuperscript{(43)} kontrastéiert. Hie féiert un enger Kette e dresséierte
\VUgras{Bier} \textsuperscript{(44)}.

%%K t09-tit-7 sub
\VUpk{10.2pt}{40}{Bierger besteigen.}
\VUsaut{3.0mm}

%%K t09-c1-16-1 p
16. --- Bal all Dag ginn ech mat engem \VUgras{Guide} \textsuperscript{(48)}
\VUgras{Klammen} üben, an ech sinn ganz houfreg, wann ech mat engem \VUgras{Sak}
\textsuperscript{(47)} um Réck an dem „\VUgras{Alpenstock}“ an der Hand, wéi e
richtegen \VUgras{Tourist} \textsuperscript{(46)}, d'\VUgras{Fielsen}
\textsuperscript{(45)} stiermen.

%%K t09-c1-17-1 p
17. --- Vun uewen um Bierg gesinn d'Awunner vun der Ebene net méi grouss aus wéi
Seechomessen; eng \VUgras{Schofsherd} \textsuperscript{(50)}, déi op enger
\VUgras{Weed} \textsuperscript{(49)} weid, gesäit aus wéi e Kannerspillsaach. E
\VUgras{Pin} \textsuperscript{(51)} oder och en \VUgras{Holzhaischen}
\textsuperscript{(52)} gesäit kaum aus wéi e Fleck am Gréngen.

%%K t09-c1-18-1 p
18. --- Bei deenen Ausfluch begéine mer dacks um \VUgras{Wee} \textsuperscript{(53)}
e \VUgras{Gamsjeeër} \textsuperscript{(54)}, deen d'Déier op der Schëller dréit, dat
hie grad geschoss huet.

%%K t09-c1-19-1 p
19. --- Ech hunn a mäin Herbarium en ganz bemierkenswäerten Zweig \VUgras{Faarn}
\textsuperscript{(55)} an e schéint rosa Zweigelchen \VUgras{Heed}
\textsuperscript{(57)} geluecht, déi ech beim Agank vun enger klenger \VUgras{Grott}
\textsuperscript{(56)} bei mengem leschte Spadséiergang geplakt hunn.

%%K t09-tit-8 sub
\VUsaut{3.0mm}
\VUcentre{12.0pt}{0}{\textit{Zweet Zeen.}}
\VUsaut{3.0mm}

%%K t09-tit-9 sub
\VUpk{11.4pt}{40}{De Bësch. --- D'Juegd.}
\VUsaut{3.0mm}

%%K t09-c2-01-1 p
1. --- Gëschter hunn ech en herrlechen Dag am Bësch verbruecht. De Fläiss, deen ënner
den Déieren an den \VUgras{Insekten} \textsuperscript{(5)} herrscht, vun deenen en
bewunnt gëtt, huet mech ganz interesséiert.

%%K t09-c2-01-2 p
D'\VUgras{Spann} \textsuperscript{(1)}, déi hiert \VUgras{Netz} \textsuperscript{(2)}
tëscht zwee Zweige spënnt, fir déi \VUgras{Méck} \textsuperscript{(3)} ze fänken, déi
laanschtflitt, d'\VUgras{Schleek} \textsuperscript{(4)}, déi mat Méi hiert Haus
dréit, den Hirschkäfer, deen no sengem Fang sicht, déi fläisseg Seechomes, ganz
äifreg bei hirem \VUgras{Seechomesenhaascht} \textsuperscript{(6)}, --- all déi Kleng
sinn gaangen, komm, hunn geschafft mat aussergewéinlechem Fläiss.

%%K t09-c2-02-1 p
2. --- Eng \VUgras{Schlaang} \textsuperscript{(7)}, déi ech op den éischte Bléck fir
eng \VUgras{Kräizotter} gehalen hunn, déi awer näischt anescht war wéi eng einfach
\VUgras{Ringelnatter,} hat sech ënner engem Bamstamm verstoppt an huet vun Zäit zu
Zäit hiert dënnt Käppchen erausgestreckt, dat ëmmer prett schéngt ze bäissen. Virun
mir ass eng déck \VUgras{Nacktschleek} \textsuperscript{(8)} schwéier gekroch, nodeems
si eng vun de schmackhaften \VUgras{Äerdbieren} \textsuperscript{(11)} verschléckt
hat, déi do am Iwwerfloss wuessen.

%%K t09-c2-03-1 p
3. --- De Buedem war mat \VUgras{Bëschblummen} bedeckt:
\VUgras{Schlësselblummen} \textsuperscript{(9)}, \VUgras{Immergréng}
\textsuperscript{(10)}, a vill aner Planzen, déi ech net kannt hunn, hunn tëscht
de \VUgras{Grasbüschelen} \textsuperscript{(12)} gebléit.

%%K t09-c2-04-1 p
4. --- An där Géigend ass d'\VUgras{Trüffel} bal esou heefeg wéi de
\VUgras{Champignon} \textsuperscript{(14)}, an e \VUgras{Ferkel}
\textsuperscript{(13)}, dat engem Fierschter gehéiert, huet gierdeg gesicht, ob et
him gelénge géif, e puer dovun ze fannen.

%%K t09-tit-10 sub
\VUsaut{3.0mm}
\VUpk{10.2pt}{40}{Wilddéiferei.}
\VUsaut{3.0mm}

%%K t09-c2-05-1 p
5. --- Ech war beim \VUgras{Enn} vun enger Zeen dobäi, déi mech ganz ameséiert huet.
Mat Hëllef vun der Nues vu sengem \VUgras{niddregen Hond} \textsuperscript{(15)} huet
de \VUgras{Fierschter} \textsuperscript{(16)} grad e \VUgras{Wilddéif}
\textsuperscript{(22)} iwwerrascht, an dem Ablack, wou dëse sech prett gemaach huet,
d'\VUgras{Wëld} \textsuperscript{(17)} matzehuelen, dat hie geschoss hat. Well hie
léiwer säi Fang opginn huet, wéi sech festhuele ze loossen, ass de Wilddéif fortgelaf,
an en \VUgras{Hues} \textsuperscript{(18)}, eng \VUgras{Hirschkou}
\textsuperscript{(19)}, e \VUgras{Réi} \textsuperscript{(20)} an e
\VUgras{Wëllschwäin} \textsuperscript{(21)} sinn um Gras leie bliwwen an hunn de
Fierschter frou gemaach.

%%K t09-c2-06-1 p
6. --- D'Villfalt vun de Beem, déi de Bësch huet, ass erstaunlech. D'\VUgras{Buchen}
\textsuperscript{(23)} an d'\VUgras{Bierken} \textsuperscript{(26)}, d'\VUgras{Pinnen,}
déi d'\VUgras{Harz} \textsuperscript{(27)} ginn, d'\VUgras{Eechen}
\textsuperscript{(29)} a vill anerer mëschen hir Zweiger duercherneen.

%%K t09-c2-06-2 p
Den \VUgras{Efeu} \textsuperscript{(24)}, deen um Stamm vun den héije Beem hänkt,
faarft de \VUgras{Bësch} \textsuperscript{(25)} mat engem glänzege Gréng, dat sech
agreabel mat der méi hell a bleech grénger Faarf vum \VUgras{Moos}
\textsuperscript{(31)} verbënnt, an deem de \VUgras{Spiecht} \textsuperscript{(30)}
Insekte sicht.

%%K t09-tit-11 sub
\VUsaut{3.0mm}
\VUpk{10.2pt}{40}{Bëschlandschaft.}
\VUsaut{3.0mm}

%%K t09-c2-07-1 p
7. --- Déi al Eechen hunn mech bezaubert. Hir déck verdréinte \VUgras{Wuerzelen}
\textsuperscript{(32)}, hallef aus dem Buedem erausgewuess, ginn hinnen e prächtegt
Bild vu Kraaft a Vigeur, an ech froe mech, wéi de \VUgras{Besëtzer}
\textsuperscript{(35)} sech dozou entschléisse kann, si ëmzehaen an se der
\VUgras{Säg} \textsuperscript{(34)} ze iwwerginn, mat där den \VUgras{Holzhaker}
\textsuperscript{(33)} schafft. Déi aarm \VUgras{Stämm} \textsuperscript{(36)}, vun
hirer \VUgras{Rinn} \textsuperscript{(37)} beraabt oder mat der \VUgras{Aaxt}
\textsuperscript{(38)} gespléckt, déi den \VUgras{Deegléiner} \textsuperscript{(39)}
hieft, maache mech traureg.

%%K t09-c2-08-1 p
8. --- Nieft drun huet en ale \VUgras{Kuelebrenner} \textsuperscript{(41)} mat senger
\VUgras{Schëpp} en \VUgras{Haufen} \textsuperscript{(42)} Kuelen opgehäift, an eng al
Fra huet e \VUgras{Bierdel} \textsuperscript{(40)} \VUgras{doudegt Holz} aus de
Zweiger vun den ëmgehaene Beem matgedroen.

%%K t09-tit-12 sub
\VUsaut{3.0mm}
\VUpk{10.2pt}{40}{D'Juegd.}
\VUsaut{3.0mm}

%%K t09-c2-09-1 p
9. --- Ech wollt e puer Ablécker am \VUgras{Fierschterhaus} \textsuperscript{(46)}
raschten goen, mä wéi ech beim \VUgras{Weier} \textsuperscript{(43)} ukoum, op deem e
schéine \VUgras{Schwan} \textsuperscript{(45)} schwëmmt, hunn ech gemierkt, datt eng
rezent \VUgras{Iwwerschwemmung} \textsuperscript{(44)} de Wee an e richtege
\VUgras{Sumpf} verwandelt hat. Ech hu missen ëmdréinen, a wéi ech laanscht d'\VUgras{Zeder}
\textsuperscript{(49)}, d'\VUgras{Pappelen} \textsuperscript{(48)} an den
\VUgras{Ulm} \textsuperscript{(47)} gaangen sinn, déi um Rand vum Bësch stinn, hunn
ech aus engem \VUgras{Jonghecken} \textsuperscript{(54)} e wonnerschéinen
\VUgras{Hirsch} \textsuperscript{(50)} erauskomme gesinn, dee vun enger hartnäckeger
\VUgras{Juegdhondsmeute} \textsuperscript{(51)} verfollegt gouf, déi och dat
\VUgras{Horn} \textsuperscript{(53)} vun de \VUgras{Reidjeeër} \textsuperscript{(52)}
ugefeiert huet. Dat aarmt Déier war ganz erschöpft. Et ass e puer Schrëtt vu mir
stierwend ëmgefall.

%%K t09-c2-10-1 p
10. --- De Jong vum Fierschter, deen de Kaméidi vun der \VUgras{Hetzjuegd}
ugezunn hat, ass aus dem Haus komm, a mir zwee si nees an de Bësch gaangen. Hien hat
eng \VUgras{Ekster} \textsuperscript{(56)} op engem Zweig vun enger aler Eech gesinn.
Hien huet gemengt, hiert Nascht wier wuel an de \VUgras{Blieder}
\textsuperscript{(58)} verstoppt, an hie wollt d'Nascht huelen.

%%K t09-c2-10-2 p
Flénk wéi e \VUgras{Kaweechelchen} \textsuperscript{(61)} ass hie vun
\VUgras{Zweig} \textsuperscript{(55)} zu Zweig geklommen, mä hien huet näischt fonnt
an huet et nëmme fäerdeg bruecht, eng al \VUgras{Kréi} \textsuperscript{(60)}
opzejagen, déi op engem \VUgras{Aascht} \textsuperscript{(57)} souz.

\VUsaut{4.0mm}
\VUfilet{22mm}
